\section{Double-Shell Tank Annulus Constraints and Requirements}
\label{sec:dst-constraints}
\begingroup
\setlength{\emergencystretch}{2em}

The vacuum crawler must operate within the access, clearance, and cleaning constraints of the DST annulus. These constraints establish requirements for the complete system, including the chassis, extraction inlet, vacuum hose, cameras, and control connections. The dimensions specified below form the preliminary design basis supplied for this project; their applicability to the selected deployment tank must be confirmed against tank drawings and available access locations.

Deployment requires the crawler to pass through a nominal 12-inch-diameter riser and transition onto the annulus floor. The access arrangement considered in this project consists of four risers spaced at approximately $90^{\circ}$ intervals around the tank. This arrangement provides alternative entry points for reaching cleaning areas in support of subsequent inspection. It does not, by itself, establish that every part of the annulus is accessible. The crawler's deployment envelope must include protruding components and hose connections, with clearance for insertion and retrieval. Its lifting arrangement must also support controlled lowering, orientation, and recovery through the same opening.

\begin{figure}[H]
\centering
\includegraphics[width=0.8\textwidth]{Tank_layout_placeholder.png}
\caption{Tank layout placeholder}
\label{fig:Tank layoutplaceholder}
\end{figure}


Once deployed, the crawler must travel and maneuver within a nominal 32-inch annular spacing in the project geometry. This value is a project design input, distinct from the approximate 30-inch annulus description in the background discussion. The usable floor corridor is further constrained by the curved transition between the tank wall and bottom, together with installed equipment. Consequently, maneuverability must be evaluated using the actual cross-sectional profile at the crawler's operating height, rather than treating the nominal spacing as a rectangular passage. The chassis, inlet, and attached hose must remain within this available clearance during forward travel, reversing, and turns. Hose routing must permit these motions without snagging or destabilizing the crawler.



The mobility requirement also includes traversal of 0.75-inch-high obstacles representative of the pipes identified in the project description as protective conduits for thermocouple instrumentation used for temperature monitoring. The obstacle height is a preliminary test target; the pipe outside diameter, supports, and installed height must be confirmed before establishing the final test geometry. The crawler must cross these obstacles with its vacuum hose attached, retain sufficient traction and stability to continue driving, and avoid damaging the instrumentation or tank surfaces. This requirement applies to the combined effects of wheel geometry, ground clearance, mass distribution, and hose loading.

\begin{figure}[H]
\centering
\includegraphics[width=0.8\textwidth]{annulus_dimension.png}
\caption{Curvature Inside The Annulus floor and Pipe}
\label{fig:Curvature Inside The Annulus floor and pipe}
\end{figure}

The primary cleaning function is to vacuum loose debris from the annulus floor and convey it through the extraction hose into a collection system. The target materials include the corrosion-related debris and fine powder residue described in the preceding section. The required debris-size range is \textbf{[minimum size] to [maximum size], [units]}, pending confirmation with the project mentor. The inlet and conveying path must accommodate the agreed range without persistent blockage, while the collection arrangement must retain the recovered material. Cleaning performance will be assessed using a defined debris mixture and hose configuration. The required recovery fraction, cleaning rate, and collection capacity remain to be established before performance testing.

Remote operation requires continuous video feedback and onboard lighting sufficient to identify obstacles, monitor the extraction inlet, and guide deployment and retrieval. Camera placement must provide a usable view during both forward and reverse motion. The operator must be able to command vehicle motion and vacuum operation and identify a loss of communication or video. A proposed control requirement is that loss of the motion-command link stops the drive system and prevents automatic restart when communication returns. Remote visual observation is also identified as an operational need in earlier Hanford annulus inspection requirements \citep{pardini2001DSTRequirements}.

\begin{figure}[H]
\centering
\includegraphics[width=0.8\textwidth]{CTF_Picture.png}
\caption{Cold Testing Facility in Richland, WA}
\label{fig:Cold Testing Facility in Richland, WA}
\end{figure}




\par
\endgroup





